% =====================================================================
%  §1 - Introduction
%
% =====================================================================

\section{Introduction}
\label{sec:intro}

Scientific and clinical decision-making depends on evidence from
heterogeneous sources, especially the primary literature. Today,
this evidence is largely consumed by humans and encoded indirectly
into annotation databases and local rules.

Existing resources such as HL7 Fast Healthcare Interoperability
Resources (FHIR) Evidence and EvidenceVariable~\cite{fhir_evidence},
the Evidence and Conclusion Ontology (ECO)~\cite{eco2019}, the
Scientific Evidence and Provenance Information Ontology
(SEPIO)~\cite{sepio}, and the Global Alliance for Genomics and
Health (GA4GH) Genomic Knowledge Standards~\cite{ga4gh_va}
advance interoperability in clinical research. However, they do not
comprehensively represent the evidence types and workflows of the
basic sciences. They focus mainly on clinical trials, high-level
assertions, and variant representations, with limited support for the
fine-grained, domain-specific evidence from basic and pre-clinical
research needed for automated reasoning and AI-ready infrastructure.

Each scientific domain demands a tailored approach for structuring
and communicating evidence. For example, the genetic evidence
discussed in this paper is generated in diverse contexts, from
experimental model systems to large-scale population studies, that
are not fully supported by current standards.
To properly use evidence derived from
scientific literature in clinical workflows, two principal challenges
require standardisation: \emph{Relevance and Specificity}, that is,
how precisely a genetic finding maps to an individual patient,
disease context, or biological hypothesis; and \emph{Confidence},
that is, whether a given piece of evidence can be robustly and
safely used, ensuring defensible compliance, interpretability, and
reproducibility.

A well-designed schema for genetic evidence helps make literature
review a \emph{semantic-parsing} problem: mapping natural-language
claims to a shared, source-anchored representation. Large language
models are increasingly capable of this kind of parsing, but
dependable automation still requires a carefully structured target
language. The present work offers one such increment. Building on decades of work by
the biomedical-ontology community, it takes a small step toward literature
review that is materially assisted by automated parsing.

\paragraph{Contributions.} This paper makes the following contributions:
\begin{itemize}
    \item A domain-analysis-driven semantic model for genetic evidence, with
    core classes (\ScientificEvidence, \EvidenceVariable,
    \EvidenceAssertion), genetic specialisations, a compact dimensional
    vocabulary, and a validation schema in the W3C Shapes Constraint
    Language (SHACL)~\cite{w3c-shacl}.
    \item An application of the model to six genetic-evidence papers spanning
        human, population, model-organism, and polygenic-score genetics,
        supported by extraction tooling.
    \footnote{\url{https://github.com/ForomePlatform/genetic-evidence-model}}
    \item A human--AI collaborative annotation workflow, packaged as executable
        AI skills for annotation and review that apply the model under curator
        supervision.
    \item A structural alignment of the core classes with HL7 FHIR Evidence,
        including a credibility representation that decomposes certainty into
        SEPIO-anchored facets. The dimensions are designed to align with ECO,
        SEPIO, IAO, and OBI; Supplementary Note~SN5 provides a partial
        term-level mapping.
    \item A curated term-level UMLS crosswalk of the dimensional vocabulary, in
        which every value is mapped with a typed relation or recorded as an
        argued gap. The crosswalk was produced with the open-source,
        model-agnostic \texttt{GEM Mapping Studio}, a standalone curation tool
        distributed as a Python package that supports the full workflow, from
        defining dimension axes in an empty workspace to adjudicating
        value-level mappings. The workflow follows a reproducible protocol with
        pre-registered search sweeps, structured rejection rationales, and a
        public decision log in the repository.
\end{itemize}

The remainder of the paper reviews related standards
(\Cref{sec:background}), defines the semantic model
(\Cref{sec:model}), and illustrates it with two worked examples
(\Cref{sec:examples}). It then describes the annotation workflow and
pilot statistics (\Cref{sec:workflow}), discusses findings, candidate
extensions, limitations, and future work (\Cref{sec:discussion}), lists
the released artefacts (\Cref{sec:availability}), and concludes
(\Cref{sec:conclusion}).
